\documentclass[11pt]{article}
\usepackage[margin=0.82in]{geometry}
\usepackage{amsmath,amssymb,mathtools}
\usepackage{booktabs,tabularx}
\usepackage[dvipsnames]{xcolor}
\usepackage[colorlinks=true,linkcolor=NavyBlue,urlcolor=BrickRed,hypertexnames=false]{hyperref}
\usepackage{microtype}
\usepackage{fancyhdr}
\usepackage{enumitem}

\definecolor{ink}{HTML}{142B37}
\definecolor{accent}{HTML}{2B6681}
\definecolor{pale}{HTML}{EDF6FA}
\definecolor{passgreen}{HTML}{16794A}
\color{ink}
\pagestyle{fancy}
\fancyhf{}
\lhead{Binary polyhedral groups}
\rhead{Proof and verification}
\cfoot{\thepage}
\setlength{\headheight}{14pt}
\newcommand{\Sfun}{\mathcal S}
\newcommand{\Sym}{\operatorname{Sym}}
\newcommand{\PASS}{\textcolor{passgreen}{\textbf{PASS}}}
\newcommand{\summarybox}[1]{\begin{center}\fcolorbox{accent}{pale}{\parbox{0.91\linewidth}{#1}}\end{center}}

\title{\textbf{Binary Polyhedral Groups in $SU(2)$}\\[3pt]
\large Paired lifts, parity cancellation, and direct matrix verification}
\author{Computational verification note}
\date{17 July 2026}

\begin{document}
\maketitle

\summarybox{\textbf{Outcome.} Explicit $2\times2$ complex matrix groups of
orders $24$, $48$, and $120$ were reconstructed for $2T$, $2O$, and $2I$.
For each group, all 67 multivariate coefficients through total degree eight
agree with the proposed paired-lift formula.  Closure, unitarity, determinant
one, the central element $-I$, and the odd-degree filter all pass.}

\section{The master coefficient formula}

For a finite matrix group $G\leq GL(V)$,
\begin{equation}
\Sfun_{G,V}(\mathbf t)=\frac1{|G|}\sum_{g\in G}
\prod_{a\geq1}\det(I-t_ag)^{-1}.                            \tag{1}
\end{equation}
The identity
\[
\det(I-tg)^{-1}=\sum_{d\geq0}\chi_{\Sym^dV}(g)t^d
\]
shows that
\begin{equation}
[t_1^{\tau_1}\cdots t_\ell^{\tau_\ell}]\Sfun_{G,V}
=\dim\left(\bigotimes_a\Sym^{\tau_a}V\right)^G.          \tag{2}
\end{equation}
Indeed, the coefficient is the trace of the Reynolds projector on the tensor
product.  Thus (1) and (2) are exact formal-power-series identities.

\section{Proof of the binary-lift formula}

Let $\widetilde\Gamma\leq SU(2)$ be the inverse image of a finite rotation
group $\Gamma\leq SO(3)$.  If $R\in\Gamma$ rotates through angle $\theta$,
its two lifts $q$ and $-q$ have spectra
\[
(e^{i\theta/2},e^{-i\theta/2}),\qquad
(-e^{i\theta/2},-e^{-i\theta/2}).
\]
Define
\begin{align*}
\Psi_\theta^+(\mathbf t)&=\prod_a
(1-2\cos(\theta/2)t_a+t_a^2)^{-1},\\
\Psi_\theta^-(\mathbf t)&=\prod_a
(1+2\cos(\theta/2)t_a+t_a^2)^{-1},\\
B_\theta&=\tfrac12(\Psi_\theta^++\Psi_\theta^-).
\end{align*}
Pair the two terms in (1) over each $R\in\Gamma$.  Because
$|\widetilde\Gamma|=2|\Gamma|$, this gives
\begin{equation}
\Sfun_{\widetilde\Gamma,\mathbb C^2}
=\frac1{|\Gamma|}\sum_\theta N_\Gamma(\theta)B_\theta.    \tag{3}
\end{equation}
This proves the proposed formula without requiring a separate conjugacy table
for the binary group.

Substituting the rotation counts yields
\begin{align}
\Sfun_{2T}&=\frac1{12}(B_0+8B_{2\pi/3}+3B_\pi),            \tag{4}\\
\Sfun_{2O}&=\frac1{24}(B_0+6B_{\pi/2}+8B_{2\pi/3}+9B_\pi),\tag{5}\\
\Sfun_{2I}&=\frac1{60}(B_0+15B_\pi+20B_{2\pi/3}
+12B_{2\pi/5}+12B_{4\pi/5}).                              \tag{6}
\end{align}

\section{The central parity filter}

The matrix $-I$ belongs to each binary group.  On
$W_\tau=\bigotimes_a\Sym^{\tau_a}(\mathbb C^2)$ it acts as
$(-1)^{|\tau|}$.  If $|\tau|$ is odd and $v$ is invariant, then
$v=(-I)v=-v$, so $v=0$.  Therefore
\begin{equation}
[m_\tau]\Sfun_{2T}=[m_\tau]\Sfun_{2O}=[m_\tau]\Sfun_{2I}=0
\quad\text{for odd }|\tau|.                                \tag{7}
\end{equation}
The same cancellation is visible termwise in the average defining
$B_\theta$.

\section{Independent matrix construction}

The parent rotation matrices are constructed geometrically: signed
permutation matrices give the octahedral group; the subgroup preserving the
four vertices $(1,1,1),(1,-1,-1),(-1,1,-1),(-1,-1,1)$ gives the tetrahedral
group; and compatible vertex-frame maps of the icosahedron give the
icosahedral group.

For each rotation matrix, a standard trace-and-diagonal algorithm recovers one
unit quaternion $q=w+xi+yj+zk$.  It is mapped to
\begin{equation}
\rho(q)=\begin{pmatrix}
w+iz&y+ix\\-y+ix&w-iz
\end{pmatrix}\in SU(2).                                    \tag{8}
\end{equation}
Both $\rho(q)$ and $-\rho(q)$ are retained.  Since the sign choice for an
individual lift is arbitrary but both signs are present, the resulting set is
the full preimage and is independent of those choices.

This construction does not use (4)--(6).  The implementation verifies all
pairwise products against the reconstructed set, as well as unitarity,
determinant one, and the presence of $-I$.

\section{Verification protocol and results}

For every partition $\tau$ through total degree eight, the direct route
recursively computes complete homogeneous characters from all explicit
matrices and averages their products.  The independent route uses only the
parent angle counts and paired half-angle spectra in (3).  All odd rows are
also checked separately against zero.

Finally, explicit matrices on $V\otimes V$ and $\Sym^2(V)$ are averaged into
Reynolds projectors.  The alternating determinant tensor gives a one-dimensional
invariant in $V\otimes V$, while $\Sym^2(V)$ has no invariant vector.

\begin{center}
\begin{tabularx}{\linewidth}{@{}Xrrrrc@{}}
\toprule
Group & order & dimension & degree & checks & result\\
\midrule
$2T$ & 24 & 2 & $0$--$8$ & 67 & \PASS\\
$2O$ & 48 & 2 & $0$--$8$ & 67 & \PASS\\
$2I$ & 120 & 2 & $0$--$8$ & 67 & \PASS\\
\bottomrule
\end{tabularx}
\end{center}

All three sets are closed under multiplication.  The largest observed
unitarity error is $3.17\times10^{-16}$ and the largest determinant error is
$2.27\times10^{-16}$.  Every odd-degree coefficient is exactly zero after
integer validation.

For each group, the $4\times4$ Reynolds projector on $V\otimes V$ has rank and
trace $1$; its largest idempotence error is $1.04\times10^{-16}$.  The
symmetric-square projector has rank $0$.  Consequently the computation
recovers the proposed low-degree distinction
\[
[m_{(2)}]=0,\qquad [m_{(1,1)}]=1.
\]

\section{Reproducibility and scope}

Run from the repository root:
{\scriptsize
\begin{verbatim}
.venv/bin/python -m for_this_guy.matrix_group_formula_verification.binary_polyhedral.verification
\end{verbatim}
}
Open \texttt{marimo\_notebooks/binary\_polyhedral.py} with marimo for interactive
degree selection and complete row tables.

The finite checks validate the matrix lift, angle convention, class counts,
central signs, and coefficient code in reasonable low dimensions.  The
all-degree conclusion is the formal proof (1)--(3); it does not depend on the
chosen cutoff or floating-point experiment.

\end{document}
